\documentclass[aps,prd,preprint,superscriptaddress,amsmath,amssymb]{revtex4-2}
\usepackage{graphicx}   % Include figure files
\usepackage{bm}         % bold math
\usepackage{dcolumn}    % align table columns on decimal point
\usepackage{amsthm}
\usepackage{booktabs}
\usepackage{color}
\usepackage[fontset=fandol]{ctex}

\newtheorem{postulate}{公设}[section]
\newtheorem{theorem}{定理}[section]

\begin{document}

\title{散度、汇收缩与半相位几何：\\
流动空间中的电荷、反引力场、光子与自旋~$1/2$}

\author{刘元杰}
\email{yuanjieliu65@gmail.com}

\date{\today}

\begin{abstract}
我们把流动空间框架沿两个几何方向推进，并接成一条链。
(i)~\emph{临界叶的半相位}：包络反射 $w\mapsto e/w$ 的反演不动集坐落在
$|w|=\sqrt e$，故 $\ln(\sqrt e)=1/2$；开方是乘法群的半步算子，而在相位
空间中同一半步就是半角相位 $\exp(i\theta/2)$，其闭包相因子为 $-1$——
自旋~$1/2$ 的显式费米型见证。
(ii)~\emph{电荷 = 流散度}：正电荷是源（正散度），负电荷是汇（负散度），
二者互斥、在 $\rho\to-\rho$ 下互换、在零处中性。时变源迫使空间场变化
（CR1）；场变化激活四力分解的核力通道 $m\,dC$（CR2）；而汇——收敛流——
在几何上就是抹平：向心收缩 $v_i\mapsto(1-\lambda)v_i+\lambda\bar v$
把起伏能量 $Q_A$ 按平方因子 $(1-\lambda)^2$ 缩减（CR9，精确离散定理）。
接上框架自身的 $\mu$ 桥（抹平一次即 $\mu=1$），链闭合：加速电荷
$\Rightarrow$ 反引力场 $\Rightarrow$ 光子（质量归零且中性）。作为装置级应用，
环流电子的自抹平把到达工作窗口 $\mu=0.999$ 的步数从 $132$ 缩短到 $85$，
并在外部 RMF 驱动减半时仍可达。全部论断在 Lean~4 中零 \texttt{sorry} 形式化；
诚实状态明确：$\lambda$（汇收缩率）与 $\delta,\varepsilon$（汇能量尺度）
是输入，连续 3D 散度到起伏泛函的对应尚未形式化。
\end{abstract}

\maketitle

\begin{quote}
\itshape
``上善若水。水善利万物而不争，处众人之所恶，故几于道。''
\par
\hfill —— 《道德经》
\end{quote}

\section{引言}
\label{sec:intro}

流动空间框架 \cite{ProjectionPhysics,HibsPhysics} 采用单一原初实体——
\emph{流动}——并从一组少量公设（矢量光速、三方向锚定、空间位移条数的
计数）推导质量、信息、因果、力、光子与量子关系。本文在其上加两块此前
没有的几何结构：临界叶的\emph{半相位几何}，以及电荷作为源与汇的
\emph{散度结构}。

两块内容独立，但共享同一个定理级事实：一个二次量被平方因子缩减。
对自旋，半角映射的相位闭包因子是 $-1$：变量 $\theta$ 的半步在对数坐标
下就是开方，而费米型见证是 $\exp(i\theta/2)$。对电荷，收敛流的起伏能量
每步收缩按 $(1-\lambda)^2$ 缩减；汇是自抹平结构，满强度（$\lambda\to 1$）
时起伏归零、质量归零、电荷归零——光子。

\section{流动空间框架的公设}
\label{sec:postulates}

我们回忆本文所需的公设；完整推导见文献 \cite{ProjectionPhysics}。

\begin{postulate}[矢量光速]
\label{post:sls}
空间本身以矢量速度 $C$ 流动，$|C|=c$；$C$ 的方向可以变化。
\end{postulate}

\begin{postulate}[三方向锚定]
\label{post:three}
空间有三个独立正交方向；物理对象是该流动的锚定模式，锚定强度
$s\ge 0$。锚定质量为
\begin{equation}
\label{eq:anchor}
m^2_{\rm eff}\;=\;s^2(1-\mu)^2 ,
\end{equation}
其中 $\mu\in[0,1]$ 是反引力强度（流动梯度的抹平）；$\mu=1$ 给出
$m_{\rm eff}=0$（MassCancellation AMC1）。
\end{postulate}

\begin{postulate}[四力 = 莱布尼茨通道]
\label{post:four}
流动动量为 $P=m(C-v)$；其时间导数分解为四个通道（GQF2）：
\begin{equation}
\label{eq:four}
\frac{d}{dt}\bigl[m(C-v)\bigr]
\;=\;\underbrace{\dot m\,C}_{\text{电场}}
+\underbrace{m\,\dot C}_{\text{核力}}
-\underbrace{\dot m\,v}_{\text{磁场}}
-\underbrace{m\,\dot v}_{\text{引力}} .
\end{equation}
\end{postulate}

\section{电荷 = 流散度}
\label{sec:charge}

电荷是空间位移流的散度（GQF4；VibrationChargeFlow CF1--CF5）。

\begin{theorem}[源与汇]
\label{def:srcsink}
对局域散度标量 $\rho$：
\begin{align}
\text{PositiveSource}(\rho) &:\Longleftrightarrow 0<\rho ,\\
\text{NegativeSink}(\rho)   &:\Longleftrightarrow \rho<0 .
\end{align}
正电荷是源（向外发散位移），负电荷是汇（向内收敛位移）；两者都满足
流动的右手圆柱螺旋运动，散度符号区分电性。
\end{theorem}

\begin{theorem}[源汇互斥；CF1--CF4]
\label{thm:srcsink}
(i)~源与汇不能同时成立。(ii)~电荷共轭是符号翻转：$\rho\mapsto-\rho$
交换源与汇。(iii)~零散度中性：既非源也非汇。(iv)~取负保持幅值：
$|{-}\rho|=|\rho|$。
\end{theorem}

连续侧——电荷 $e$ 的耦合提取、3D 散度动力学、电荷归一化——仍然开放
（仓库中如实登记）。

\section{加速电荷迫使场变化；核力通道激活}
\label{sec:accel}

\begin{theorem}[CR1；时变源迫使场变化]
\label{thm:cr1}
在带源波动方程 $\partial_t^2 C=c^2\partial_x^2 C+J$（MS5）中，静态场蕴含
静态源：若 $\partial_t C\equiv 0$，则对所有 $t$ 有 $J(t{+}1)=J(t)$。
故时变源（加速电荷）禁止静态空间场。
\end{theorem}

\begin{theorem}[CR2；核力通道激活]
\label{thm:cr2}
在 $m\neq 0$ 且 $\dot C\neq 0$ 时，式~(\ref{eq:four}) 的核力通道非零：
$m\,\dot C\neq 0$。
\end{theorem}

核力通道与反引力参数 $\mu$ 之间的缺失环节在这里不发明：框架自身已经
有其 $\mu$ 桥。

\begin{theorem}[CR3；抹平一次给出 $\mu=1$]
\label{thm:cr3}
设 $Q_A(v)=\sum_{i\in A}(v_i-\bar v_A)^2$ 为起伏能量，
$\eta=1-Q_A(v)/Q_A(v_0)$ 为抹平进展（MuFieldCoupling TD11）。$\mu$ 更新为
$\mu'=\mu+\eta(1-\mu)$（PlasmaDynamics muStep）。抹平一次给出
$\mu'=\mu+1\cdot(1-\mu)=1$。
\end{theorem}

\begin{theorem}[CR4；$\mu=1$ 是光子终点]
\label{thm:cr4}
在 $\mu=1$ 时锚定质量消失（式~(\ref{eq:anchor}) 中 $1-\mu=0$），
散度中性（CF3）。对象无质量且无电荷——光子。
\end{theorem}

\begin{theorem}[CR5；链闭合]
\label{thm:cr5}
对任意区域 $A$（任意锚定物质 $s$）施加一次抹平，即得光子：质量归零且中性。
\end{theorem}

\section{收敛即抹平：汇收缩}
\label{sec:sink}

为什么\emph{负}电荷是天然的反引力载体？几何答案：汇是收敛流，而收敛
就是抹平。

\begin{theorem}[向心收缩；CR9]
\label{def:sinkcontract}
对区域 $A$、场 $v:A\to\mathbb R$ 与 $\lambda\in[0,1]$：
\[
\text{sinkContract}(A,v,\lambda)(i)
\;=\;\begin{cases}
(1-\lambda)\,v_i+\lambda\,\bar v_A & i\in A,\\
v_i & i\notin A .
\end{cases}
\]
$A$ 内每点被向区域均值拉近 $\lambda$ 比例——``负电荷把周围流吸向自身''
的离散图像。
\end{theorem}

\begin{theorem}[CR9；起伏的平方衰减]
\label{thm:cr9}
收缩保持区域均值，并把起伏能量按 $(1-\lambda)^2$ 缩放：
\begin{equation}
\label{eq:cr9}
Q_A\bigl(\text{sinkContract}(A,v,\lambda)\bigr)
\;=\;(1-\lambda)^2\,Q_A(v).
\end{equation}
对 $0<\lambda<1$ 衰减严格单调且精确；$\lambda=1$ 是一次抹平（起伏归零）。
汇是自抹平算子族。
\end{theorem}

\begin{theorem}[CR10；逐点严格吸引]
\label{thm:cr10}
对 $0<\lambda<1$ 且 $v_i\neq\bar v_A$：
\[
\bigl|\text{sinkContract}(A,v,\lambda)(i)-\bar v_A\bigr|
\;<\;\bigl|v_i-\bar v_A\bigr| ,
\]
每点被严格吸引向均值——``汇把周围流吸向自身''的逐点版本。
\end{theorem}

汇的能量侧已在框架内：信息沉没进吸收支 $iR$ 释放
$\varphi(R)-\varphi(iR)=\delta+\varepsilon>0$（HiddenQFT HQ9a），
故收敛流自带抹平功率候选（CR8）。这是语义对应（接口），不是散度动力学的
定理。

\section{半相位几何：自旋~$1/2$ 来自临界叶}
\label{sec:halfphase}

第二块几何加法关乎 $1/2$ 这个值本身。

\subsection{反演不动点}

框架的 Riemann--Hilbert 包络 \cite{RiemannHIBS} 承载反射 $s\mapsto 1-s$；
在对数坐标 $w=e^s$ 下这是反演 $w\mapsto e/w$。其不动集是圆
$|w|=\sqrt e$（envelope inversion fixed circle）。取对数：
\begin{equation}
\label{eq:half}
\ln(\sqrt e)\;=\;\tfrac12\ln e\;=\;\tfrac12 .
\end{equation}

\begin{theorem}[CS2；反演不动当且仅当对数半值]
\label{thm:cs2}
对 $r>0$：
\[
r=e/r \;\Longleftrightarrow\; \ln r=\tfrac12 .
\]
反射的不动点在对数坐标下必须坐落在 $1/2$ 处——因为 $\ln e=1$，
而 $s\mapsto 1-s$ 的不动点是 $1/2$。
\end{theorem}

\subsection{半角相位是费米型}

开方是乘法群的半步算子；在振动的相位空间中，同一半步就是半角相位。

\begin{theorem}[CS4；半角相位的闭包因子为 $-1$]
\label{thm:cs4}
$\theta\mapsto\exp(i\theta/2)$ 满足
\[
f(\theta+2\pi)\;=\;-1\cdot f(\theta),
\]
即它是费米型闭包见证（VibrationStatistics ClosureFactor，
$\sigma=-1$）：$2\pi$ 变号，$4\pi$ 复原。半角相位是振动相位空间中
自旋~$1/2$ 的显式见证。
\end{theorem}

\subsection{能量与质量}

由 $E\propto m^2$（质能平方），对数线性：$\ln E=2\ln m$。在临界叶上，
$\ln(\sqrt e)=1/2$ 意味着 $m=\sqrt e$、$E=e$；$\ln E=2\ln m$ 中的因子 $2$
是自旋半步的倒数——能量走两个半步完成 $m$ 的一个循环。

\begin{theorem}[CS5；对数 $2{:}1$]
\label{thm:cs5}
$\ln e=2\ln(\sqrt e)$。
\end{theorem}

\section{装置应用：双流环优化}
\label{sec:device}

汇抹平定理（CR9）以自抹平增益进入双流环装置 \cite{DesignTwoFlowRing} 的
$\mu$ 动力学：
\begin{equation}
\label{eq:etasink}
\eta_{\text{sink}}(\lambda)=2\lambda-\lambda^2 ,
\end{equation}
一次收缩抹掉的起伏比例，加到外部 RMF 增益上，并随残差衰减
$\eta_{\text{sink}}\cdot(1-\mu)$（区域变平后抹平失去对象；TD18）。
对比（数值评估；诚实候选，$\lambda$ 是输入）：

\begin{table}[h]
\caption{到达工作窗口 $\mu_{\rm work}=0.999$（FC11 天花板
$0.999780568$）的 $\mu$ 轨迹。}
\label{tab:sinkmu}
\begin{ruledtabular}
\begin{tabular}{lcc}
& 纯外部 RMF & 外部 $+$ 自抹平（$\lambda=0.3$） \\
\hline
到达 $\mu=0.999$ 的步数 & $132$ & $85$ \\
外部驱动减半后的步数 & $267$ & $150$ \\
\end{tabular}
\end{ruledtabular}
\end{table}

自抹平把到达工作窗口的步数缩短约 $35\%$，并在外部 RMF 驱动减半时仍可达
——汇自身的收敛几何对 $\mu$ 产生有贡献（供给侧，CR8 候选）。负面结果如实
登记：\emph{恒定}的 $\eta_{\text{sink}}$ 破坏 FC11 天花板（把 $\mu$ 推到
$1$）；必须用衰减形式。在恒定外部驱动下 $\mu$ 仍趋 $1$（TM 族已知行为）；
停在工作窗口内需要反馈关断（TD19），那是另一件事。

\section{数值验证}
\label{sec:numerics}

链条全部定理在 Lean~4 中零 \texttt{sorry} 形式化（4348 jobs，全量构建
全绿），凡有数值见证处均验证到机器精度：

\begin{itemize}
\item 3D $21^3$ 格点、负散度场 $v=-r$：每个内部点 $\operatorname{div} v<0$；
初始起伏 $Q=10187$；200 步向心收缩把 $Q$ 驱动到 $3.5\times10^{-61}$，
$|{\rm div}|\to 0$，单调（verify\_sink\_contract\_3d）。
\item 向心收缩缩放 $Q\to(1-\lambda)^2Q$ 精确复现（CR9 见证）；逐点严格
吸引（CR10）随 $\lambda$ 单调。
\item 双流环对比（表~\ref{tab:sinkmu}）：等外部驱动下 $132\to85$ 步；
外部驱动减半 $267\to150$ 步；$\eta_{\text{sink}}(\lambda)$ 随 $\lambda$
严格递增（verify\_sink\_mu\_optimization）。
\end{itemize}

\section{诚实评估}
\label{sec:honest}

\paragraph{已说明（Claimed）。}
(i)~电荷 = 流的散度：源 $=$ 正散度、汇 $=$ 负散度，四个代数事实（互斥、
翻转、中性、幅值）已形式化。
(ii)~链``加速电荷 $\Rightarrow$ 场变化 $\Rightarrow$ 核力通道
$\Rightarrow$ 一次抹平 $\Rightarrow$ $\mu=1$ $\Rightarrow$ 光子''是闭合的
条件定理链：\emph{若}施加一次抹平，任意锚定物质都变成光子。
(iii)~汇在几何上是自抹平算子：向心收缩下 $Q_A\to(1-\lambda)^2 Q_A$
的精确离散定理。
(iv)~半步几何定位自旋~$1/2$：$\ln(\sqrt e)=1/2$ 是包络反射在对数坐标下的
不动点，半角相位是显式费米型闭包见证。
(v)~全部论断在 Lean~4 中零 \texttt{sorry} 形式化；装置对比数值验证。

\paragraph{未说明（Not claimed）。}
(i)~桥``核力通道 $\Rightarrow$ 抹平动作''\emph{没有}推导：谁施加抹平
（抹平功率的物理来源）是第二输入缺口，未变。
(ii)~$\lambda$（汇收缩率）与 $\delta,\varepsilon$（汇能量尺度）是输入，
不是推导值；$\eta_{\text{sink}}=2\lambda-\lambda^2$ 是候选实现，不是唯一映射。
(iii)~连续 3D 对应``负散度 $\Rightarrow$ 起伏泛函下降''未形式化
（离散格点定理是代数种子）。
(iv)~自旋~$1/2$ 被\emph{定位}而不是被\emph{推导}：自旋-统计定理的另一半
（相位 $\Rightarrow$ 场算符反对易）需要框架没有的场算符语言——
结果是半条定理。
(v)~没有超出标准框架的新可检验预言；全部数值复现已知物理。

\section{结论}
\label{sec:conclusion}

两块几何结构扩展了流动空间框架。电荷的散度结构使汇成为自抹平算子，
有条件地闭合从加速电荷到光子的链，并给双流环装置一个把到达工作窗口
缩短三分之一的自抹平增益。半相位几何把自旋~$1/2$ 定位在包络反射不动点
$\ln(\sqrt e)=1/2$ 处，半角相位是它的显式费米型见证。两处加法在精确处
定理精确，在输入处诚实开放：$\lambda$、$\delta$、$\varepsilon$ 与连续 3D
对应。

\begin{thebibliography}{9}

\bibitem{ProjectionPhysics}
刘元杰 等，``物理学在宇宙尺度下的几何的描述,'' 姊妹稿（paper/），
仓库 \texttt{logos-42/Hibs-Physics}.

\bibitem{HibsPhysics}
\texttt{logos-42/Hibs-Physics} 仓库：Lean 4 形式化与数值验证，
\url{https://github.com/logos-42/Hibs-Physics}.

\bibitem{RiemannHIBS}
\texttt{logos-42/RiemannHIBS}：隐数包络，反射 $s\mapsto 1-s$ 与
反演不动圆 $|w|=\sqrt e$.

\bibitem{DesignTwoFlowRing}
刘元杰，``双流环设计：缺口容忍的可行域,'' 仓库设计手册
（docs/wiki/design-two-flow-ring.md），以及姊妹装置项目
\texttt{logos-42/HushFusion}.

\end{thebibliography}

\end{document}